\documentclass[11pt,oneside,openany]{book}

\usepackage[T1]{fontenc}
\usepackage{lmodern}
\usepackage[letterpaper,margin=1.10in,includeheadfoot]{geometry}
\usepackage{amsmath,amssymb,amsthm,mathtools}
\usepackage{microtype}
\usepackage{xcolor}
\usepackage{titlesec}
\usepackage{fancyhdr}
\usepackage{needspace}
\usepackage{hyperref}

\ifdefined\pdftrailerid
  \pdftrailerid{MathGenome ALG-03 core supplement}
\fi

\definecolor{ink}{HTML}{172923}
\definecolor{deepgreen}{HTML}{2B5242}
\definecolor{rulegreen}{HTML}{8A9F92}
\definecolor{palegreen}{HTML}{F0F4EF}
\hypersetup{
  colorlinks=true,linkcolor=deepgreen,urlcolor=deepgreen,citecolor=deepgreen,
  pdftitle={Rings, Modules, and Associative Algebras},
  pdfsubject={Jacobson density, central simple algebras, full corners, and quiver representations},
  pdfauthor={MathGenome},bookmarksnumbered=true
}
\setcounter{tocdepth}{1}
\setcounter{secnumdepth}{2}
\setlength{\headheight}{15pt}
\setlength{\emergencystretch}{2em}
\titleformat{\chapter}[display]
  {\normalfont\color{ink}}{\color{deepgreen}\large\scshape Chapter \thechapter}{9pt}
  {\huge\bfseries}
\titlespacing*{\chapter}{0pt}{4pt}{24pt}
\titleformat{\section}{\normalfont\Large\bfseries\color{ink}}{\thesection}{0.7em}{}
\titleformat{\subsection}{\normalfont\large\bfseries\color{ink}}{\thesubsection}{0.7em}{}

\pagestyle{fancy}
\fancyhf{}
\fancyhead[L]{\small\sffamily\color{deepgreen}MathGenome}
\fancyhead[R]{\small\nouppercase{\leftmark}}
\fancyfoot[C]{\small\thepage}
\renewcommand{\headrulewidth}{0.35pt}
\renewcommand{\chaptermark}[1]{\markboth{\thechapter.\ #1}{}}
\fancypagestyle{plain}{%
  \fancyhf{}\fancyfoot[C]{\small\thepage}\renewcommand{\headrulewidth}{0pt}}

\newtheorem{theorem}{Theorem}[chapter]
\newtheorem{lemma}[theorem]{Lemma}
\newtheorem{proposition}[theorem]{Proposition}
\newtheorem{corollary}[theorem]{Corollary}
\theoremstyle{definition}
\newtheorem{definition}[theorem]{Definition}
\newtheorem{example}[theorem]{Example}
\theoremstyle{remark}
\newtheorem{remark}[theorem]{Remark}
\allowdisplaybreaks[1]

\begin{document}
\frontmatter
\hypersetup{pageanchor=false}
\begin{titlepage}
  \thispagestyle{empty}
  \thispdfpagelabel{Cover}
  \begin{flushleft}
    \vspace*{0.65in}
    {\large\scshape\color{deepgreen}MathGenome\par}
    \vspace{0.32in}
    {\color{rulegreen}\rule{0.94\textwidth}{0.65pt}\par}
    \vspace{0.45in}
    {\color{ink}\fontsize{34}{39}\selectfont\bfseries
       Rings, Modules, and\par Associative Algebras\par}
    \vspace{0.30in}
    {\color{deepgreen}\Large
       Density, Brauer Groups, Full Corners,\\[3pt]
       and Quiver Representations\par}
    \vfill
    {\large Supplementary lecture notes\par}
    \vspace{0.12in}
    {\normalsize ALG-03 \enspace\textperiodcentered\enspace MathGenome v0.22.0 reference baseline\par}
    \vspace*{0.48in}
  \end{flushleft}
\end{titlepage}
\hypersetup{pageanchor=true}
{\small\tableofcontents}
\chapter*{Scope and reading guide}
\addcontentsline{toc}{chapter}{Scope and reading guide}
\markboth{Scope and reading guide}{}

These notes supplement \emph{Algebra Core} and \emph{Abstract Algebra}
\cite{p103,d07} with Jacobson density, the basic theory of central simple
algebras, Morita equivalence for full idempotents, and representations
of quivers with relations.

Rings and algebra maps are unital. Modules are left modules in the
density and quiver chapters and right modules in the full-corner chapter.
A \(k\)-algebra has a central structure map. Endomorphisms multiply by
composition and the rightmost path in a product acts first.

{\setlength{\fboxsep}{0.8em}\setlength{\fboxrule}{0.35pt}%
\noindent\fcolorbox{rulegreen}{palegreen}{%
\begin{minipage}{\dimexpr\linewidth-2\fboxsep-2\fboxrule\relax}
\textbf{Result keys.} \texttt{RM03-C03.01} belongs to module
\texttt{ALG-03-C03}. The source map records each result's PDF and TeX
location and its proof dependencies. Keys beginning with \texttt{p103:}
refer to \emph{Algebra Core}.
\end{minipage}}\par}

\section*{The six core modules}

\paragraph{C01: Rings, modules, and exact sequences.}
\emph{Algebra Core}, Chapter~1, proves the quotient, exactness, change
of scalars, splitting, and diagram results. \emph{Abstract Algebra},
Chapters~9--11 and~16, treats factorization, Smith form, and PID modules.

\paragraph{C02: Tensor, Hom, projective, injective, and flat modules.}
See \emph{Algebra Core}, Chapters~1--2, for balanced tensor, tensor--Hom,
exactness, dual bases, and the projective, injective, and flat criteria.
The right progenerator result is \texttt{p103:ALG-03-C02-S05}.

\paragraph{C03: Chain conditions, radicals, and semisimplicity.}
\emph{Algebra Core}, Chapters~2--3, covers finite length, Schur,
semisimplicity, the radical, Nakayama, and Artin--Wedderburn.
Chapter~\ref{ch:rm03-density} adds density and its finite-dimensional quotient.

\paragraph{C04: Central simple algebras.}
\emph{Abstract Algebra}, Chapter~23, provides the division ring inputs.
Chapter~\ref{ch:rm03-central-simple} proves the splitting, conjugacy,
centralizer, and Brauer group results. See also the descent comparison
in \emph{Linear Algebraic Groups}, Result~F05.

\paragraph{C05: Morita equivalence and idempotents.}
\emph{Algebra Core}, Chapter~3, proves the progenerator theorem and
the matrix case. Chapter~\ref{ch:rm03-morita} gives the full-corner criterion
and functors. See \emph{Algebraic K-Theory}, Result~A10 \cite{ktheory}.

\paragraph{C06: Representations of associative algebras.}
For group algebras, Schur, and Maschke, see \emph{Algebra Core},
Chapter~8 (\texttt{p103:ALG-07-C01-R01--R03}).
Chapter~\ref{ch:rm03-quivers} proves the quiver comparisons and examples.


\mainmatter
\chapter{Jacobson density}
\label{ch:rm03-density}

A faithful simple module makes its ring act by linear maps over a
division ring. The density theorem describes exactly what this action
can do on any finite list of vectors. The inputs from \emph{Algebra Core}
are Schur's lemma and the splitting of submodules of a semisimple module
\cite[Ch.~3, ALG-03-C03-R03]{p103}. A version over a field appears in
\cite[Thm.~3.3]{mitdensity}.

\section{Finite interpolation on a simple module}

\begin{theorem}[RM03-C03.01. Jacobson density]
\label{thm:rm03-density}
Let \(R\) be a unital ring and \(M\) a nonzero simple left \(R\)-module.
Put \(D=\operatorname{End}_R(M)^{\mathrm{op}}\), and make \(M\) a right
\(D\)-space by \(m\cdot f^{\mathrm{op}}=f(m)\). For every finite
right \(D\)-linearly independent list \(x_1,\ldots,x_n\) in \(M\) and
every \(y_1,\ldots,y_n\in M\), there is \(r\in R\) such that
\[
   rx_i=y_i\qquad (1\leq i\leq n).
\]
\end{theorem}
\begin{proof}
Schur's lemma makes \(D\) a division ring. Multiplication in the
opposite ring reverses composition, so the displayed evaluation is a
right action. Every \(f\in\operatorname{End}_R(M)\) commutes with left
multiplication by elements of \(R\). Thus the \(R\)-action induces a
ring map \(R\to\operatorname{End}_D(M)\) with composition as product.

For an empty list take \(r=0\). Otherwise give \(M^n\) the diagonal
left \(R\)-action and write \(x=(x_1,\ldots,x_n)\). Suppose \(Rx\) is
proper. The module \(M^n\) is semisimple, so \(Rx\) has a nonzero
complement \(U\) and there is an \(R\)-linear projection \(p:M^n\to U\).
Some coordinate projection \(\pi_j:M^n\to M\) has nonzero restriction
to \(U\), since a nonzero vector has a nonzero coordinate. Composing
\(p\) with that restriction gives a nonzero map \(h:M^n\to M\) which
vanishes on \(Rx\). For the coordinate inclusions \(\iota_i:M\to M^n\),
write \(f_i=h\iota_i\in\operatorname{End}_R(M)\). The identity
\(h(z_1,\ldots,z_n)=\sum_i f_i(z_i)\) gives
\[
      0=h(x)=\sum_i x_i\cdot f_i^{\mathrm{op}}.
\]
Right \(D\)-linear independence forces each \(f_i=0\), contrary to
\(h\ne 0\). Consequently \(Rx=M^n\), so a single \(r\) sends \(x\)
to \((y_1,\ldots,y_n)\).
\end{proof}

\begin{corollary}[RM03-C03.02. Density in the finite topology and the finite-dimensional quotient]
\label{cor:rm03-density-finite}
In the setting of Theorem~\ref{thm:rm03-density}, the image of
\(R\to\operatorname{End}_D(M)\) is dense for the topology of agreement
on finite subsets of \(M\). If \(\dim_D M\) is finite, the induced map
\[
      R/\operatorname{Ann}_R(M)\longrightarrow\operatorname{End}_D(M)
\]
is an isomorphism of rings.
\end{corollary}
\begin{proof}
For \(T\in\operatorname{End}_D(M)\) and a finite subset \(S\subseteq M\),
choose a right \(D\)-basis of its span. Apply the theorem to that basis
and its images under \(T\). Both \(T\) and the action of the resulting
\(r\) are right \(D\)-linear, so they agree on the span and hence on
\(S\). The sets of maps agreeing with \(T\) on such finite subsets
form a neighbourhood basis at \(T\). If \(M\) is finite-dimensional,
agreement on a full basis gives \(T\) itself. The kernel of the action
is precisely the two-sided ideal of elements annihilating \(M\).
\end{proof}

\begin{example}[RM03-C03.03. A proper dense image]
\label{ex:rm03-density-proper}
Let \(V\) have countably infinite basis \(v_0,v_1,\ldots\) over a
field \(k\), and let \(R=k\,1+F(V)\subseteq\operatorname{End}_k(V)\),
where \(F(V)\) is the ideal of finite-rank maps. Then \(V\) is a
faithful simple left \(R\)-module with commutant \(k\,1\), and the
action of \(R\) on \(V\) has proper dense image.
\end{example}
\begin{proof}
The action is faithful by construction. For \(v\ne0\) and any \(w\),
choose a linear functional \(\ell\) with \(\ell(v)=1\). The rank-at-most-one
map \(u\mapsto\ell(u)w\) sends \(v\) to \(w\), proving simplicity.
An \(R\)-endomorphism \(T\) of \(V\) is \(k\)-linear because the scalar
maps lie in \(R\). Commuting with all coordinate projections forces
\(T(v_i)=c_i v_i\) for \(c_i\in k\). Commuting with the matrix unit
sending \(v_i\) to \(v_j\) gives \(c_i=c_j\). Thus the commutant is
exactly the scalars. Density now follows from
Corollary~\ref{cor:rm03-density-finite}. The shift \(S(v_i)=v_{i+1}\)
lies outside \(R\), since the vectors \(v_{i+1}-c v_i\) are linearly
independent for every \(c\in k\). Indeed, the coefficient of the
highest indexed \(v_{i+1}\) in any finite relation is the coefficient
of the corresponding summand. Hence \(S-c\,1\) has infinite rank for
every \(c\).
\end{proof}

\chapter{Central simple algebras and the Brauer group}
\label{ch:rm03-central-simple}

Throughout this chapter \(k\) is a field. Algebras and algebra maps are
unital, and a simple algebra is nonzero and has no nonzero proper
two-sided ideal. For a field extension \(L/k\), write \(A_L=A\otimes_k L\).
We use the Artin--Wedderburn theorem from
\cite[ALG-03-C03-R04]{p103} and the finite-dimensional case of density
from Corollary~\ref{cor:rm03-density-finite}. For comparison with the
classical development, see \cite[Sections~11.3--11.8]{stackscsa}.
The finite and real division-ring inputs are the established results in
\cite[K07.23.finite\_division\_ring and K07.23.frobenius]{d07}.

\section{Structure and extension of scalars}

\begin{definition}[RM03-C04.01. Central simple algebras]
\label{def:rm03-csa}
A \emph{central simple \(k\)-algebra} is a finite-dimensional simple
\(k\)-algebra \(A\) whose center is \(k1_A\). A \emph{central division
algebra} over \(k\) is a finite-dimensional division algebra with this
same center condition. A central simple \(k\)-algebra is \emph{split}
if it is isomorphic to \(M_n(k)\) for some \(n\geq 1\). An extension
\(L/k\) \emph{splits} \(A\) if \(A_L\) is split over \(L\).
\end{definition}

\begin{proposition}[RM03-C04.02. Matrix and division presentations]
\label{prop:rm03-csa-matrix}
Every central simple \(k\)-algebra \(A\) admits a presentation
\(A\simeq M_r(D)\), where \(r\geq 1\) and \(D\) is a central division
algebra over \(k\). Both \(r\) and the \(k\)-algebra isomorphism class
of \(D\) are uniquely determined by \(A\).
\end{proposition}
\begin{proof}
A finite-dimensional simple algebra is left Artinian, so the simple
case of Artin--Wedderburn gives \(A\simeq M_r(D)\). Commuting with
the matrix units forces any central matrix to have the form \(dI_r\),
and commuting with coefficient matrices forces \(d\in Z(D)\).
Thus \(Z(D)=k\), with the identification induced by the original
\(k\)-action on \(A\). For a simple left \(A\)-module \(S\) the matrix
presentation gives \(S\simeq D^r\), with endomorphism ring \(D^{\mathrm{op}}\)
via right multiplication by \(D\). This identifies \(D\) as a
\(k\)-algebra since every module endomorphism commutes with the
given scalars. Finally \(r^2\dim_k D=\dim_k A\) determines \(r\).
\end{proof}

The same simple-Artinian structure will be used for a finite-dimensional
simple \(k\)-algebra \(T\) with arbitrary center. Write \(T=M_s(E)\)
for its Artin--Wedderburn presentation. A finite nonzero left
\(T\)-module \(V\) is a direct sum of \(m\geq 1\) copies of the
unique simple column module \(E^s\). In particular the
\(k\)-dimension of \(V\) determines this multiplicity. Its
endomorphism ring is \(M_m(E^{\mathrm{op}})\), by writing an
endomorphism in blocks between the \(m\) copies of \(E^s\).
Consequently \(\operatorname{End}_T(V)\) is simple, and
\begin{equation}
\label{eq:rm03-simple-end-dimension}
\dim_k T\;\dim_k\operatorname{End}_T(V)=(\dim_k V)^2.
\end{equation}
Indeed the two factors on the left are \(s^2\dim_k E\) and
\(m^2\dim_k E\), and \(\dim_k V=sm\dim_k E\).

\begin{proposition}[RM03-C04.03. The enveloping algebra of a central simple algebra]
\label{prop:rm03-enveloping}
If \(A\) is central simple over \(k\), its left and right multiplication
map is an isomorphism of \(k\)-algebras
\[
 A\otimes_k A^{\mathrm{op}}\longrightarrow \operatorname{End}_k(A),
 \qquad a\otimes b^{\mathrm{op}}\longmapsto (x\mapsto axb).
\]
\end{proposition}
\begin{proof}
The displayed map is an algebra homomorphism because composition
reverses the order of the right multipliers. Under this action the
left \(A\otimes A^{\mathrm{op}}\)-submodules of \(A\) are exactly its
two-sided ideals, so \(A\) is a simple module. Its endomorphisms as
such a module are determined by the image \(c\) of \(1\). Compatibility
with left and right multiplication gives \(ac=ca\) for every \(a\),
and each such \(c\) defines an endomorphism. The resulting endomorphism
division algebra is \(k\). Corollary~\ref{cor:rm03-density-finite} therefore
gives surjectivity onto \(\operatorname{End}_k(A)\). Both \(k\)-dimensions
are \((\dim_k A)^2\), which gives injectivity.
\end{proof}

\begin{proposition}[RM03-C04.04. Tensor product and base change]
\label{prop:rm03-csa-tensor}
Let \(A\) be central simple over \(k\) and \(B\) any simple
\(k\)-algebra. Then \(A\otimes_k B\) is simple, with center
\(1\otimes Z(B)\). In particular the tensor product of two central
simple \(k\)-algebras is central simple. Also \(A_L\) is central
simple over \(L\) for every field extension \(L/k\).
\end{proposition}
\begin{proof}
Let \(I\) be a nonzero two-sided ideal in \(A\otimes_k B\). Express
a nonzero \(u\in I\) as \(\sum_i a_i\otimes b_i\) in a basis
\((a_i)\) of \(A\), and choose \(j\) with \(b_j\ne 0\).
The ideal is preserved by left and right multiplication from
\(A\otimes 1\), hence is preserved by \(T\otimes\operatorname{id}_B\)
for every \(k\)-linear \(T:A\to A\) by
Proposition~\ref{prop:rm03-enveloping}. Take \(T(a_j)=1\) and
\(T(a_i)=0\) for \(i\ne j\). This gives \(1\otimes b_j\in I\).
Because \(Bb_jB\) generates the unit ideal of \(B\), left and
right multiplication by \(1\otimes B\) then gives \(1\otimes 1\in I\).

An element commuting with \(A\otimes 1\) can be written
\(\sum_j x_j\otimes y_j\) for \(k\)-linearly independent
\(y_j\in B\). Each \(x_j\) commutes with \(A\), so the element
belongs to \(1\otimes B\). An element of that subalgebra is central
exactly when it commutes with \(1\otimes B\), giving the center formula.
Applying it to \(B=L\), and noting
\(\dim_L A_L=\dim_k A\), gives the base change statement.
\end{proof}

\section{Separable splitting fields}

\begin{theorem}[RM03-C04.05. Degree and finite Galois splitting]
\label{thm:rm03-csa-galois-split}
For \(A\) central simple over \(k\), \(\dim_k A\) is a square.
Write \(\deg A=\sqrt{\dim_k A}\). There is a finite Galois extension
\(K/k\) with \(A_K\simeq M_{\deg A}(K)\).
More precisely, if \(A\simeq M_r(D)\) and \(d=\deg D\), then \(D\)
contains a separable field extension \(E/k\) of degree \(d\) which
splits \(D\). Thus \(\deg A=rd\).
\end{theorem}
\begin{proof}
Over an algebraic closure \(\overline{k}\), base change and
Proposition~\ref{prop:rm03-csa-matrix} give
\(A_{\overline{k}}\simeq M_n(D')\).
Every element \(x\) of \(D'\) satisfies a polynomial over
\(\overline{k}\) by finite dimensionality. Factoring the polynomial
into linear factors and using absence of zero divisors shows
\(x\in\overline{k}\). This proves \(A_{\overline{k}}\simeq
M_n(\overline{k})\) and \(\dim_k A=n^2\). The same argument for \(D\)
gives \(D_{\overline{k}}\simeq M_d(\overline{k})\) and
\(n=rd\).

It suffices to split \(D\), since matrix algebras commute with
extension of scalars. If \(k\) is finite, then \(D\) is a finite
division ring and is commutative by
\cite[K07.23.finite\_division\_ring]{d07}. Its center condition forces
\(D=k\), so \(E=k\) works.

Assume \(k\) is infinite. Fix a \(k\)-basis \(u_1,\ldots,u_{d^2}\)
of \(D\). Through the splitting over \(\overline{k}\), the matrix
corresponding to \(\sum_i X_i u_i\) has a characteristic polynomial
whose discriminant is a polynomial in the \(X_i\) over
\(\overline{k}\). This discriminant is nonzero as a polynomial:
the basis spans \(M_d(\overline{k})\) after extension, so choose
coefficients giving a diagonal matrix with \(d\) distinct entries.
Any nonzero polynomial over \(\overline{k}\) in finitely many
variables has a nonzero value at some tuple in \(k\). To see this,
view it as a polynomial in its last variable, choose the other
variables by induction so that at least one coefficient survives,
and then avoid the finitely many roots of the remaining polynomial.
Thus some \(x\in D\) corresponds to a matrix with \(d\) distinct
eigenvalues over \(\overline{k}\).

The minimal polynomial of \(x\) has the same degree after extension:
linear independence of the initial powers \(1,x,\ldots,x^{j-1}\)
is preserved, and a monic relation over \(k\) remains a relation
over \(\overline{k}\). Uniqueness of the monic relation of smallest
degree then identifies the two minimal polynomials. For the chosen
matrix this polynomial is the product of the \(d\) distinct linear
factors. Hence \(E=k[x]\) has dimension \(d\) and is separable over \(k\).
The algebra \(E\) is a field since it is a finite-dimensional
commutative domain inside \(D\).

Regard \(D\) as a right \(E\)-space of dimension \(d\). The map
\[
 D\otimes_k E\longrightarrow \operatorname{End}_E(D),\qquad
 a\otimes e\longmapsto (z\mapsto aze)
\]
is an \(E\)-algebra homomorphism: the right multipliers belong to
the commutative field \(E\). It is unital and its domain is simple
by Proposition~\ref{prop:rm03-csa-tensor}, so it is injective.
Both sides have \(E\)-dimension \(d^2\), giving a splitting over
\(E\). A finite Galois closure \(K\) of \(E/k\) then splits both \(D\)
and \(A\).
\end{proof}

\begin{corollary}[RM03-C04.06. Index and degrees of splitting extensions]
\label{cor:rm03-csa-index}
Let \(A\simeq M_r(D)\) be central simple and set
\(\operatorname{ind}(A)=\deg D\). Every finite splitting field
\(L/k\) of \(A\) has \(\operatorname{ind}(A)\mid [L:k]\).
A separable splitting field of degree \(\operatorname{ind}(A)\)
exists.
\end{corollary}
\begin{proof}
Write \(d=\deg D\). Since \(A_L=M_r(D_L)\) and \(D_L\) is central
simple by Proposition~\ref{prop:rm03-csa-tensor}, the uniqueness
of the division presentation implies that \(D_L\) is split
whenever \(A_L\) is split. Then \(D_L\simeq M_d(L)\), and
restricting its natural left action on \(L^d\) to \(D\) gives a
finite left \(D\)-space. Its \(k\)-dimension \(d[L:k]\) must be a
multiple of \(\dim_k D=d^2\). The last assertion is the field \(E\)
of Theorem~\ref{thm:rm03-csa-galois-split}.
\end{proof}

\section{Conjugacy and commutants}

\begin{theorem}[RM03-C04.07. Skolem--Noether]
\label{thm:rm03-skolem-noether}
Let \(A\) be central simple over \(k\), and let \(B\) be a
simple \(k\)-algebra. Two unital \(k\)-algebra embeddings
\(f,g:B\hookrightarrow A\) are conjugate by an element of \(A^\times\).
In particular every \(k\)-algebra automorphism of \(A\) is inner.
\end{theorem}
\begin{proof}
The embedding makes \(B\) finite-dimensional over \(k\), and
\(T=A\otimes_k B^{\mathrm{op}}\) is a finite-dimensional simple
\(k\)-algebra by Proposition~\ref{prop:rm03-csa-tensor}.
Give the \(k\)-space \(A\) two left \(T\)-module structures by
\[
 (a\otimes b^{\mathrm{op}})\cdot_f x=axf(b),\qquad
 (a\otimes b^{\mathrm{op}})\cdot_g x=axg(b).
\]
The use of \(B^{\mathrm{op}}\) compensates for the order of right
multiplication. The two modules have equal \(k\)-dimension, so the
simple-Artinian module description preceding
Proposition~\ref{prop:rm03-enveloping} gives a \(T\)-module
isomorphism \(F:(A,\cdot_f)\to(A,\cdot_g)\).
Because \(F\) is left \(A\)-linear, \(F(x)=xu\) for \(u=F(1)\).
Its inverse is also right multiplication, which shows that
\(u\in A^\times\). Intertwining the two actions of \(b^{\mathrm{op}}\)
on \(1\) gives \(f(b)u=ug(b)\). The case \(B=A\) and \(g\) the
identity gives the automorphism assertion.
\end{proof}

\begin{theorem}[RM03-C04.08. Centralizer and double centralizer]
\label{thm:rm03-centralizer}
Let \(B\) be a unital simple \(k\)-subalgebra of a central simple
\(k\)-algebra \(A\). Its centralizer
\[
 C_A(B)=\{\,x\in A:xb=bx\ \text{for every }b\in B\,\}=C
\]
is simple, \(\dim_k B\cdot\dim_k C=\dim_k A\), and \(C_A(C)=B\).
Moreover \(Z(C)=Z(B)\) as subalgebras of \(A\).
If \(B\) is central over \(k\), so is \(C\), and multiplication
induces an isomorphism \(B\otimes_k C\simeq A\).
\end{theorem}
\begin{proof}
Act on \(A\) by \(T=A\otimes B^{\mathrm{op}}\) using left and right
multiplication, where \(T\) is simple by
Proposition~\ref{prop:rm03-csa-tensor}. An endomorphism of this
\(T\)-module is left \(A\)-linear, hence is right multiplication
by some \(x\in A\). Right multiplication by \(x\) commutes with
right multiplication by every \(b\in B\) exactly when \(x\in C\).
Thus \(\operatorname{End}_T(A)\simeq C^{\mathrm{op}}\).
By \eqref{eq:rm03-simple-end-dimension}, \(C\) is simple and
\[
 (\dim_k A)(\dim_k B)(\dim_k C)=(\dim_k A)^2.
\]
The same dimension formula applied to the simple subalgebra
\(C\) gives \(\dim_k C_A(C)=\dim_k B\). The inclusion
\(B\subseteq C_A(C)\) therefore is an equality. Also
\(Z(C)=C\cap C_A(C)=C\cap B=Z(B)\).

If \(Z(B)=k\), multiplication is a unital homomorphism from
\(B\otimes_k C\) because the subalgebras commute. The tensor
product is simple by Proposition~\ref{prop:rm03-csa-tensor},
so the homomorphism is injective. The dimension formula gives
surjectivity.
\end{proof}

\begin{corollary}[RM03-C04.09. Maximal subfields of a central division algebra]
\label{cor:rm03-maximal-subfield}
Every maximal subfield \(E\) of a central division \(k\)-algebra
\(D\) has \([E:k]=\deg D\) and splits \(D\). The separable \(E\)
in Theorem~\ref{thm:rm03-csa-galois-split} is a maximal subfield.
\end{corollary}
\begin{proof}
For \(x\in C_D(E)\), the algebra \(E[x]\) is a commutative
finite-dimensional domain over \(E\), hence a field. Maximality
gives \(x\in E\), so \(C_D(E)=E\). Theorem~\ref{thm:rm03-centralizer}
now gives \([E:k]^2=\dim_k D\). The map using left multiplication
by \(D\) and right multiplication by \(E\) in the proof of
Theorem~\ref{thm:rm03-csa-galois-split} works for any \(E\) of
this degree. Conversely the same dimension formula forces
\(C_D(E)=E\) for every degree-\(\deg D\) subfield, which is
sufficient for maximality.
\end{proof}

\section{Reduced characteristic polynomial and norm}

\begin{proposition}[RM03-C04.10. Reduced characteristic polynomial]
\label{prop:rm03-reduced-norm}
Let \(A\) be central simple of degree \(n\) over \(k\).
For \(a\in A\), the characteristic polynomial of the image of
\(a\otimes 1\) in a matrix algebra over a splitting field belongs
to \(k[t]\) and is independent of the field and of the splitting
isomorphism. Denote it by \(\chi_a^{\mathrm{red}}(t)\) and set
\(\operatorname{Nrd}(a)=(-1)^n\chi_a^{\mathrm{red}}(0)\).
The reduced norm is a homogeneous polynomial of degree \(n\)
on \(A\), commutes with base change, is multiplicative, and
is nonzero exactly on \(A^\times\).
\end{proposition}
\begin{proof}
Choose a finite Galois splitting \(K/k\), an isomorphism
\(\phi:A_K\to M_n(K)\), and a \(k\)-basis \(u_1,\ldots,u_{n^2}\)
of \(A\). On the generic element put
\[
 P(t,\mathbf X)=
 \det\left(tI_n-\sum_i X_i\phi(u_i\otimes 1)\right).
\]
For \(\sigma\in\operatorname{Gal}(K/k)\), write \(\sigma_A\)
for the action on the coefficients of \(A_K\) and \(\sigma_M\)
for entrywise application to matrices. Then
\(\sigma_M\phi\sigma_A^{-1}\) is another \(K\)-algebra
isomorphism. By Theorem~\ref{thm:rm03-skolem-noether} the
two differ by a single conjugation on \(M_n(K)\). Hence
applying \(\sigma\) to the coefficients leaves \(P\) unchanged.
Consequently \(P\in k[t,\mathbf X]\). If \(L/k\) is another
splitting field, an algebraic closure of \(L\) contains a
\(k\)-copy of \(K\). Extending the two splitting isomorphisms
to this common field and using the same conjugacy argument
identifies their generic characteristic polynomials. These
identities also verify base change for a field extension
that does not itself split \(A\), by using a splitting field
over it and comparing there.

The constant term with sign \((-1)^n\) is the determinant
of the matrix of the generic element. Homogeneity and
multiplicativity follow from the determinant formulas after
extension to \(K\), and hence over \(k\). If \(a\) is invertible
then its matrix image is invertible. Conversely if its
matrix image is invertible, left multiplication \(x\mapsto ax\)
on \(A\) becomes invertible over \(K\). It is therefore
invertible over \(k\), giving some \(b\) with \(ab=1\).
Injectivity of left multiplication also gives \(ba=1\)
from \(a(ba-1)=0\).
\end{proof}

Over \(K\), the norm-one equation is the determinant-one
equation defining \(\operatorname{SL}_n\). Its descent is
\(\operatorname{SL}_1(A)\) in \cite[Result F05]{lag}, which
also gives the corresponding inner form \(\operatorname{PGL}_1(A)\).

\section{Brauer classes and examples}

\begin{theorem}[RM03-C04.11. Brauer equivalence and the group law]
\label{thm:rm03-brauer-group}
Two central simple \(k\)-algebras \(A,B\) are \emph{Brauer
equivalent} if \(M_p(A)\simeq M_q(B)\) for some \(p,q\geq 1\).
This is an equivalence relation, and each class has a unique
central division representative up to \(k\)-algebra isomorphism.
The classes form an abelian group \(\operatorname{Br}(k)\)
with laws
\[
 [A]+[B]=[A\otimes_k B],\qquad 0=[k],\qquad -[A]=[A^{\mathrm{op}}].
\]
The class \([A]\) is zero exactly when \(A\) is split. A field
extension \(L/k\) gives a homomorphism
\(\operatorname{Br}(k)\to\operatorname{Br}(L)\) by \([A]\mapsto[A_L]\).
\end{theorem}
\begin{proof}
Write \(A=M_r(D)\) and \(B=M_s(E)\) with division factors.
If they become isomorphic after passing to matrices, then
\(D\simeq E\) by Proposition~\ref{prop:rm03-csa-matrix}.
If \(D\simeq E\), both \(M_s(A)\) and \(M_r(B)\) are
isomorphic to \(M_{rs}(D)\). Thus equivalence means equality
of the division factors up to \(k\)-algebra isomorphism. The
class of \(k\) therefore consists exactly of split algebras.

The tensor product is central simple by
Proposition~\ref{prop:rm03-csa-tensor}. The identity
\(M_p(A)\otimes_k B\simeq M_p(A\otimes_k B)\), induced by
tensoring coefficients, makes the product independent of
the representative of \(A\). Commutativity via the tensor
flip handles the representative of \(B\). Associativity and
the identity class come from the tensor associator and unit.
The opposite algebra is central simple and
Proposition~\ref{prop:rm03-enveloping} identifies its product
with \(A\) as the split algebra \(\operatorname{End}_k(A)\).
Finally, extension of scalars commutes with the matrix and
tensor constructions and sends \(k\) to \(L\). These
identifications also give compatibility in towers of fields.
\end{proof}

\begin{example}[RM03-C04.12. The real and quaternion classes]
\label{ex:rm03-brauer-real}
There are exactly two real Brauer classes, represented by
\(\mathbb R\) and \(\mathbb H\). To see this, apply
Frobenius's real algebraic division-ring theorem
\cite[K07.23.frobenius]{d07} to the central division
representative. Finite-dimensionality guarantees
algebraicity, and \(\mathbb C\) is excluded because its
center is larger than \(\mathbb R\). The quaternion class
is nonzero by uniqueness of the division representative.
Quaternion conjugation \(q\mapsto\bar q\) reverses products,
so \(\mathbb H^{\mathrm{op}}\simeq\mathbb H\) and
\(2[\mathbb H]=0\). Thus
\(\operatorname{Br}(\mathbb R)\simeq\mathbb Z/2\mathbb Z\).

More explicitly \(\mathbb C=\mathbb R[i]\) is the degree-two
splitting subfield of \(\mathbb H\) from
Corollary~\ref{cor:rm03-maximal-subfield}, so
\(\mathbb H_{\mathbb C}\simeq M_2(\mathbb C)\).
Write \(q=a+bi+cj+dk\) in Hamilton's basis, with
\(i^2=j^2=-1\) and \(k=ij=-ji\). In a complex splitting one
can take
\[
 i\longmapsto \begin{pmatrix}\mathrm i&0\\0&-\mathrm i\end{pmatrix},
 \qquad j\longmapsto \begin{pmatrix}0&1\\-1&0\end{pmatrix},
 \quad \mathrm i^2=-1\ \text{in }\mathbb C .
\]
These matrices satisfy the quaternion relations, giving a
unital homomorphism out of \(\mathbb H_{\mathbb C}\).
It is an isomorphism by simplicity and equal dimension.
Its determinant on \(q\) is \(a^2+b^2+c^2+d^2\), recovering
the quaternion norm of \cite[K07.23.quaternion\_algebra]{d07}.
\end{example}

\begin{remark}[RM03-C04.13. Relative classes and descent]
\label{rem:rm03-relative-brauer}
For \(K/k\) finite Galois, let
\(\operatorname{Br}(K/k)=\ker(\operatorname{Br}(k)\to\operatorname{Br}(K))\).
It consists of the classes split by \(K\), by
Theorem~\ref{thm:rm03-brauer-group}. Put \(G=\operatorname{Gal}(K/k)\).
In fixed degree \(n\), the algebras split by \(K\) up to
isomorphism correspond to \(H^1(G,\operatorname{PGL}_n(K))\)
with its coefficientwise action. We first give the descent argument.

For a finite-dimensional \(K\)-space \(V\) with a semilinear
\(G\)-action \(T_\sigma\), the fixed vectors span \(V\) over \(K\).
Otherwise a nonzero \(K\)-linear functional \(\ell\) annihilates
their \(K\)-span. The vector
\(\sum_{\sigma\in G} T_\sigma(\lambda v)\) is fixed for each
\(v\in V\), \(\lambda\in K\). Applying \(\ell\) gives
\(\sum_\sigma\ell(T_\sigma v)\sigma(\lambda)=0\) for all \(\lambda\).
Distinct field automorphisms are linearly independent as
\(K\)-valued functions: in a nonzero relation with the fewest
terms, evaluate at \(a\lambda\) and subtract a multiple of the
relation at \(\lambda\), choosing \(a\) on which two of the
automorphisms differ. This eliminates one term without
eliminating the other. A one-term relation is impossible by
evaluation at \(1\). Thus \(\ell(v)=0\) for all \(v\), a
contradiction. A \(K\)-basis of fixed vectors is a \(k\)-basis of
\(V^G\), since the coefficients of a fixed vector in this basis
are fixed by \(G\). Consequently \(K\otimes_k V^G\to V\) is an
isomorphism. If \(V\) is a \(K\)-algebra and the action respects
multiplication, this is an algebra isomorphism. A \(G\)-equivariant
\(K\)-algebra map descends by restriction to the fixed subalgebras.

Given a splitting of an algebra of degree \(n\), transport its
coefficient action to \(M_n(K)\). Composing each transported
action with inverse entrywise \(\sigma\) gives a \(K\)-algebra
automorphism, inner by Theorem~\ref{thm:rm03-skolem-noether}.
Two implementing invertible matrices differ by a scalar, since
the matrices commuting with all the matrix units are the scalar
ones. Thus the semilinear action gives a projective cocycle.
Conversely a projective cocycle acts on \(M_n(K)\) by semilinear
conjugation; the scalar defects disappear in conjugation. Its
fixed algebra is central simple of degree \(n\). Its scalar
extension is the matrix algebra, so any nonzero proper ideal
would give a nonzero proper ideal there by extending a vector-space
basis. A central element extends to a scalar matrix, and its
scalar is fixed by \(G\), hence lies in \(k\). Changing the
splitting conjugates the action by a \(K\)-algebra isomorphism,
which gives just the projective coboundary relation.
Conversely that relation gives an equivariant conjugation, which
restricts to an isomorphism on fixed algebras. This proves the
fixed-degree classification. Brauer equivalence allows
enlargement by matrices.

For the descent convention where the transported semilinear
action is \(x\mapsto g_\sigma \sigma(x)g_\sigma^{-1}\),
choose lifts \(G_\sigma\in\operatorname{GL}_n(K)\) of the
projective matrices \(g_\sigma\). Then
\(c_{\sigma,\tau}=G_\sigma\sigma(G_\tau)G_{\sigma\tau}^{-1}\)
is scalar by the projective cocycle identity. Associating the
product of three lifted maps in two ways gives
\[
 c_{\sigma,\tau}c_{\sigma\tau,\nu}
 =\sigma(c_{\tau,\nu})c_{\sigma,\tau\nu}.
\]
Changing each \(G_\sigma\) by a scalar changes \(c\) by a
coboundary. Changing the splitting conjugates the projective
datum by a coboundary, which one can lift as
\(G_\sigma\mapsto H^{-1}G_\sigma\sigma(H)\); this leaves the
scalar \(c_{\sigma,\tau}\) unchanged. Enlargement by a
split matrix algebra replaces \(G_\sigma\) by \(I_m\otimes G_\sigma\)
and preserves \(c\). Thus the scalar class depends only
on the relative Brauer class. Tensoring lifts multiplies
the scalar cocycles. We have constructed the group homomorphism
\[
 \operatorname{Br}(K/k)\longrightarrow
 H^2(\operatorname{Gal}(K/k),K^\times),\qquad [A]\longmapsto[c].
\]
With the same lifts, the cocycle in
\cite[Equation~(59.61.5.1)]{stacksdescent} is \(c^{-1}\).
The homomorphism is injective. If the scalar class is zero,
rescale the \(G_\sigma\) so that \(c_{\sigma,\tau}=1\).
The operators \(v\mapsto G_\sigma\sigma(v)\) on \(K^n\) then
form a semilinear action. In a basis of invariant vectors the
action becomes entrywise \(\sigma\). The fixed endomorphisms
are the endomorphisms over \(k\) of the fixed vector space,
so the descended algebra is split.

For surjectivity take a normalized scalar \(2\)-cocycle \(c\).
On the \(K\)-space with basis \(e_\tau\), \(\tau\in G\), define
the invertible \(\sigma\)-semilinear operators
\[
 T_\sigma(e_\tau)=c_{\sigma,\tau}e_{\sigma\tau}.
\]
On \(e_\nu\), the cocycle identity gives
\(T_\sigma T_\tau=c_{\sigma,\tau}T_{\sigma\tau}\).
The operators therefore act genuinely on the full endomorphism
algebra by conjugation. The fixed algebra descends as above.
Relative to the displayed basis the \(T_\sigma\)'s are lifts
of its projective cocycle and the scalar defect is exactly \(c\).
This proves the relative isomorphism with the indicated sign.
Compare Colliot-Th\'el\`ene and Skorobogatov,
\href{https://www.imo.universite-paris-saclay.fr/~jean-louis.colliot-thelene/BGgroup_book.pdf}{Theorem~1.3.5}.

For the separable-closure statement, put \(G_k=\operatorname{Gal}(k^s/k)\)
and give \((k^s)^\times\) the discrete topology. A continuous cochain
on \(G_k^r\) has finite image and factors through a finite quotient:
choose product neighborhoods on which it is constant, cover the
compact domain by finitely many of these neighborhoods, and take a
common open normal subgroup refining all their factors. Enlarge the
finite Galois field fixed by that subgroup to contain the finitely
many cochain values. The cocycle equation, and any equality with a
coboundary, can then be checked on the finite quotient associated to
that enlarged field. This identifies continuous second cohomology
with the filtered colimit of the finite Galois groups'
second cohomology under inflation and coefficient inclusion.
The relative class maps agree with these transitions: extend the
splitting and use the same lifts inflated along restriction of
automorphisms. Every class of \(\operatorname{Br}(k)\) is split
at a finite Galois stage by Theorem~\ref{thm:rm03-csa-galois-split}.
Taking the union of relative classes therefore gives
\(\operatorname{Br}(k)\cong H^2_{\mathrm{cts}}(G_k,(k^s)^\times)\)
for this convention. Compare \cite[Theorem~59.61.6]{stacksdescent}.
\end{remark}

\chapter{Morita equivalence for full corners}
\label{ch:rm03-morita}

We use the right-module convention in this chapter. The progenerator
form of Morita equivalence and its converse are proved in
\cite[Ch.~3, ALG-03-C02-S05]{p103}. In that theorem a finitely generated
projective generator \(P_R\) gives inverse equivalences by tensoring
with \(\operatorname{Hom}_R(P,R)\) and with \(P\) on the indicated
sides. Here we identify the generator condition and the two functors
for a corner
\cite[Lec.~7]{mitmorita}.

\section{The full-idempotent criterion}

\begin{corollary}[RM03-C05.01. Morita equivalence from a full idempotent]
\label{cor:rm03-full-corner}
Let \(R\) be a unital ring and \(e\in R\) satisfy \(e^2=e\). The
right module \(eR\) is finitely generated projective, and the map
\[
   eRe\longrightarrow\operatorname{End}_R(eR),\qquad
   a\longmapsto (x\longmapsto ax)
\]
is an isomorphism of unital rings when \(eRe\) has identity \(e\).
The module \(eR\) is a generator of right \(R\)-modules if and only if
\(e\) is full, meaning that \(ReR=R\). In this case there are inverse
equivalences of right-module categories
\[
   \mathrm{Mod}\text{-}R \longrightarrow \mathrm{Mod}\text{-}eRe,
   \quad M\longmapsto Me,
   \qquad
   N\longmapsto N\otimes_{eRe}eR
\]
with the second functor directed back to \(\mathrm{Mod}\text{-}R\).
\end{corollary}
\begin{proof}
The decomposition \(R_R=eR\oplus(1-e)R\) follows by multiplying on
the left by \(e\) and \(1-e\). Therefore \(eR\) is projective and is
generated by \(e\). Every \(R\)-endomorphism \(f\) of \(eR\) has
\(f(e)=f(e)e\in eRe\) and satisfies \(f(er)=f(e)r\). Evaluation at
\(e\) is inverse to the displayed left multiplication map. Products
correspond to composition and the identity endomorphism to \(e\).

If \(e\) is full, a finite expression
\(1=\sum_{i=1}^m r_i e s_i\) gives an \(R\)-linear surjection
\[
    (eR)^m\longrightarrow R_R,\qquad
    (x_i)_{i=1}^m\longmapsto\sum_i r_i x_i,
\]
since \((e s_i)_i\) maps to \(1\). Every right module is a quotient
of a direct sum of copies of \(R_R\), so every right module is a
quotient of a direct sum of copies of \(eR\). Conversely, if \(eR\)
is a generator, take a surjection from a direct sum of copies of
\(eR\) onto \(R_R\) and a finite-support preimage of \(1\). Each
component \(g_i:eR\to R\) has \(g_i(e)=g_i(e)e\) and takes \(e s\)
to \(g_i(e)e s\in ReR\). Thus \(1\in ReR\), which proves fullness.

Evaluation at \(e\) identifies \(\operatorname{Hom}_R(eR,R)\)
with \(Re\). The inverse of evaluation sends \(b\in Re\) to
\(x\mapsto bx\) on \(eR\). Multiplication on the left by \(R\)
and precomposition by \(a\in eRe\) become the usual left \(R\)- and
right \(eRe\)-actions on \(Re\), respectively. Thus the first
functor of the progenerator theorem is \(M\otimes_R Re\).
The natural right \(eRe\)-linear map \(m\otimes b\mapsto mb\)
identifies this tensor product with \(Me\). Its inverse sends
\(u\in Me\) to \(u\otimes e\). One composite is \(ue=u\), and
the other is \(mb\otimes e=m\otimes be=m\otimes b\).
The inverse functor in the theorem is \(N\mapsto N\otimes_{eRe}eR\).
\end{proof}

Under these identifications the counit is the multiplication map
\(Me\otimes_{eRe}eR\to M\). One can read off its inverse from a full
expression \(1=\sum_i r_i e s_i\). It sends \(m\) to
\(\sum_i (m r_i e)\otimes(e s_i)\). Multiplying gives \(m\) back.
For \(m'\in Me\) and \(x\in eR\), the element \(x r_i e\) belongs to
\(eRe\). Balancing moves this element across the tensor sign and
gives the other composition
\(\sum_i m'\otimes x r_i e s_i=m'\otimes x\).

\begin{example}[RM03-C05.02. A matrix corner and a non-full corner]
\label{ex:rm03-corners}
For \(n\geq 1\), \(e_{11}\) is full in \(M_n(R)\), with corner
isomorphic to \(R\). In the product of two nonzero rings \(A\times B\),
the corner of \(e=(1,0)\) is isomorphic to \(A\), and the corner
functor \(M\mapsto Me\) kills the nonzero right module \(B\).
\end{example}
\begin{proof}
The matrix units satisfy \(\sum_i e_{i1}e_{11}e_{1i}=1_{M_n(R)}\).
The corner is the position \((1,1)\), and its product is the product
of \(R\). Thus Corollary~\ref{cor:rm03-full-corner} gives the matrix
equivalence. For the product ring the ideal generated by \((1,0)\)
is \(A\times 0\), and right multiplication by \((1,0)\) on the
module \(B\) obtained from projection to \(B\) is zero.
\end{proof}

\chapter{Representations and Path Algebras}
\label{ch:rm03-quivers}

Let \(k\) be a field. Algebras and module actions are unital, and
\(\operatorname{End}_k(V)\) is multiplied by composition. We use left
modules throughout. In particular, an arrow from \(i\) to \(j\) will
act from the vertex summand at \(i\) to the vertex summand at \(j\).

\section{Representations of an algebra}

\begin{proposition}[RM03-C06.01. Algebra representations and modules]
\label{prop:rm03-algebra-representations}
For a \(k\)-algebra \(A\) and a \(k\)-vector space \(V\), unital
\(k\)-algebra homomorphisms
\(\rho\colon A\to\operatorname{End}_k(V)\) correspond bijectively to
left \(A\)-module structures on \(V\) extending its given
\(k\)-vector-space structure. A linear map intertwines two such
representations precisely when it is \(A\)-linear.

If \(I\) is a two-sided ideal of \(A\), restriction along
\(A\to A/I\) identifies left \(A/I\)-modules with left \(A\)-modules
annihilated by \(I\), with all module maps preserved.
\end{proposition}

\begin{proof}
Set \(av=\rho(a)(v)\). The linearity of \(\rho\) and of the operators
gives the distributive laws and the given scalar action.
Multiplicativity gives \((ab)v=a(bv)\), and \(\rho(1_A)=1_V\) gives
the unit law. In the other direction \(a\) acts by a linear map
because \(a(\lambda v)=(a(\lambda 1_A))v=\lambda(av)\), and the
action laws give the required \(k\)-algebra homomorphism. The two
constructions recover the same operator at every \(a\). The
intertwining condition says exactly that \(f(av)=af(v)\).

An action obtained by restriction is zero on \(I\). Conversely, if
\(IV=0\), the formula \((a+I)v=av\) is independent of the
representative and satisfies the module laws. A map commutes with the
action of \(A/I\) exactly when it commutes with the action of all
representatives in \(A\).
\end{proof}

The group case starts with \(G\to\operatorname{GL}_k(V)\), and the
linear extension to \(k[G]\) gives the action in
Proposition~\ref{prop:rm03-algebra-representations}. The group algebra
correspondence, Schur's lemma, and Maschke's theorem with its
characteristic hypothesis are established in
\cite[Ch.~8]{p103}. In the quiver case, the generators act between
specified summands, and the multiplication of paths controls which
composites are defined.

\section{Paths, idempotents, and vertex spaces}

\begin{definition}[RM03-C06.02. Quivers and path algebras]
\label{def:rm03-path-algebra}
A finite quiver \(Q=(Q_0,Q_1,s,t)\) has a finite nonempty set
\(Q_0\) of vertices, a finite set \(Q_1\) of arrows, and source and
target maps \(s,t\colon Q_1\to Q_0\). Parallel arrows and loops are
allowed. For each vertex \(i\) there is a path \(e_i\) of length zero
with source and target \(i\). A positive-length path is written
\(p=a_\ell\cdots a_1\), where \(a_1\) is traversed first and
\(t(a_b)=s(a_{b+1})\) for \(1\leq b<\ell\).

The path algebra \(kQ\) is the vector space with basis the paths.
For basis paths \(p,q\), the product \(pq\) is their concatenation
with \(q\) traversed first when \(s(p)=t(q)\), and is zero
otherwise. Concatenation with a length-zero path at a matching
vertex leaves the other path unchanged. Multiplication extends
bilinearly to finite linear combinations.
\end{definition}

\begin{lemma}[RM03-C06.03. Multiplication and the path basis]
\label{lem:rm03-path-basis}
The multiplication makes \(kQ\) an associative unital \(k\)-algebra.
The vertices give pairwise orthogonal idempotents with sum \(1\).
The space \(e_j(kQ)e_i\) has basis the paths from \(i\) to \(j\).
The algebra \(kQ\) is finite-dimensional exactly when \(Q\) has no
oriented cycle.
\end{lemma}

\begin{proof}
For paths \(p,q,r\), both \((pq)r\) and \(p(qr)\) vanish unless
\(s(p)=t(q)\) and \(s(q)=t(r)\). When both equalities hold, each
parenthesization yields the same ordered list of arrows, starting
with \(r\). Thus multiplication is associative on the basis and
hence on \(kQ\).

The length-zero paths satisfy \(e_j e_i=0\) for \(i\ne j\) and
\(e_i^2=e_i\). For a basis path \(p\), the product
\(e_j p e_i\) is \(p\) if \(s(p)=i\) and \(t(p)=j\), and is zero
otherwise. This proves the corner basis claim. Summing the left or
the right vertex action gives the identity on every path, so
\(\sum_i e_i=1\).

The powers of an oriented cycle have distinct lengths and are
linearly independent basis elements. If there is no oriented cycle,
no path visits a vertex twice. A repeated vertex would give a
positive-length closed path, and a shortest closed subpath would
give an oriented cycle. All lengths are then at most
\(\lvert Q_0\rvert-1\). Finiteness of the arrow set gives finitely
many paths of these lengths.
\end{proof}

\begin{definition}[RM03-C06.04. Quiver representations]
\label{def:rm03-quiver-representations}
A representation \(V\) of \(Q\) consists of vector spaces \(V_i\)
at the vertices and linear maps \(V_a\colon V_{s(a)}\to V_{t(a)}\)
on the arrows. For a path \(p=a_\ell\cdots a_1\) put
\(V_p=V_{a_\ell}\circ\cdots\circ V_{a_1}\), and set
\(V_{e_i}=1_{V_i}\). A morphism \(f\colon V\to W\) is a family of
linear maps \(f_i\colon V_i\to W_i\) satisfying
\(f_{t(a)}V_a=W_a f_{s(a)}\) for each \(a\). Composition is taken
at each vertex.
\end{definition}

\begin{theorem}[RM03-C06.05. Quiver representations and path algebra modules]
\label{thm:rm03-quiver-equivalence}
The category of left \(kQ\)-modules is \(k\)-linearly equivalent to
the category of representations of \(Q\). The equivalence sends a
module \(M\) to the spaces \(e_iM\) with \(a\colon i\to j\) acting
by \(m\mapsto am\) from \(e_iM\) to \(e_jM\). Its inverse sends
\(V\) to \(\bigoplus_i V_i\). The equivalence restricts to
finite-dimensional objects and identifies module exactness with
exactness at every vertex.
\end{theorem}

\begin{proof}
Since the idempotents sum to \(1\), every \(m\in M\) is a sum of
the \(e_i m\). Applying \(e_j\) to any relation among these
summands isolates the summand at \(j\), proving that the sum is
direct. By Lemma~\ref{lem:rm03-path-basis}, \(e_j a e_i=a\) for
an arrow \(a\colon i\to j\), which gives the indicated arrow map.
A module map restricts to the summands and commutes with these
arrow maps.

Conversely, on \(\bigoplus_i V_i\) let \(p\colon i\to j\) act by
projecting to \(V_i\), applying \(V_p\), and including into
\(V_j\). The operators for composable paths multiply as the
operator of their concatenation. For noncomposable paths the
intermediate inclusion followed by projection is zero. The sum of
the vertex operators is the identity. Extending the operators
linearly therefore gives a unital \(k\)-algebra homomorphism, so
Proposition~\ref{prop:rm03-algebra-representations} gives the
module. The direct sum of the components of a morphism commutes
with the path operators, by induction on path length using the
arrow condition, and hence is a module map.

For a module, the canonical sum map \(\bigoplus_i e_iM\to M\)
commutes with the vertices and the arrows. Every basis path is a
product of these generators, so the sum map is a module
isomorphism. Starting with \(V\), each \(e_i\) becomes the
projection onto \(V_i\). Recovering those summands identifies
their arrow maps with the original \(V_a\). Both identifications
commute with morphisms by their formulas.

Each module map is block diagonal on the vertex summands, and its
kernel and image are direct sums of the component kernels and
images. Thus equality of image and kernel in a sequence is
equivalent to equality at each vertex. For arbitrary quiver
morphisms the vertex kernels are preserved by the arrow maps and
the arrows descend to the vertex quotients by the images, by the
same commutation law. This gives kernels and cokernels with the
universal properties at each vertex. Finally, the vertex set is
finite, so \(\bigoplus_i V_i\) is finite-dimensional if and only if
each \(V_i\) is finite-dimensional.
\end{proof}

\section{Relations and finite-dimensional quotients}

\begin{definition}[RM03-C06.06. Relations and admissible ideals]
\label{def:rm03-quiver-relations}
A relation from \(i\) to \(j\) is a finite linear combination
\(r=\sum_p c_p p\) of paths from \(i\) to \(j\). A representation
satisfies \(r=0\) if \(\sum_p c_p V_p\colon V_i\to V_j\) is zero.
Relations generate a two-sided ideal in \(kQ\). Every two-sided
ideal \(I\) can be described this way, since \(r\in I\) is the sum
of the components \(e_j r e_i\in I\).

Let \(J\) be the two-sided ideal generated by all the arrows.
An ideal \(I\) is admissible if there is an integer \(m\geq 2\)
such that \(J^m\subseteq I\subseteq J^2\).
\end{definition}

\begin{theorem}[RM03-C06.07. The module category with relations]
\label{thm:rm03-bound-quiver-equivalence}
Let \(I\) be generated as a two-sided ideal by a set of relations.
The equivalence in Theorem~\ref{thm:rm03-quiver-equivalence}
identifies left \(kQ/I\)-modules with quiver representations
satisfying those relations, with the same morphisms. Exactness in
this category is checked at each vertex.
\end{theorem}

\begin{proof}
Under the action on \(\bigoplus_i V_i\), the operator of a
relation \(r=\sum_p c_p p\) from \(i\) to \(j\) is the projection
to \(V_i\) followed by \(\sum_p c_p V_p\), followed by inclusion
of \(V_j\). Hence the relation acts by zero exactly when its
required component is zero. If the generating relations act by
zero, each term \(arb\), for \(a,b\in kQ\) and a generating
relation \(r\), also acts by zero. This is equivalent to
annihilation by \(I\). Proposition~\ref{prop:rm03-algebra-representations}
now gives the quotient action and all its module maps. Kernels,
images, and quotients of modules annihilated by \(I\) remain
annihilated by \(I\), so the exactness criterion follows from
Theorem~\ref{thm:rm03-quiver-equivalence}.
\end{proof}

\begin{proposition}[RM03-C06.08. Finiteness from admissibility]
\label{prop:rm03-admissible-finiteness}
For \(m\geq 1\), \(J^m\) is the span of all paths of length at
least \(m\). If \(J^m\subseteq I\), the quotient \(kQ/I\) is
finite-dimensional and is spanned by the images of the paths of
length below \(m\). If also \(I\subseteq J^2\), the images of all
length-zero and length-one paths are linearly independent.
\end{proposition}

\begin{proof}
The span of the positive-length paths is a two-sided ideal, and
every such path is a product of arrows, so it equals \(J\).
Every product of \(m\) positive-length paths vanishes or has
length at least \(m\). Conversely, a path of length at least \(m\)
can be cut into \(m\) composable positive-length paths. These two
observations identify the span with \(J^m\). There are finitely
many paths of length below \(m\), since \(Q_0\) and \(Q_1\) are
finite. The additional inclusion \(I\subseteq J^2\) says that a
linear combination of paths of length zero and one lying in \(I\)
lies in the span of paths of length two and above. Independence
of the path basis makes that combination zero.
\end{proof}

For the path convention, the module equivalence, and admissible
ideals, see \cite[Secs.~1.3--1.4]{quiver}. The two smallest
examples below also test which vertex idempotents occur on
either side of an arrow and how a relation changes its operator.

\begin{example}[RM03-C06.09. One arrow and a triangular matrix algebra]
\label{ex:rm03-one-arrow}
Take \(Q\) with vertices \(1,2\) and one arrow \(a\colon 1\to 2\).
The path basis is \((e_1,e_2,a)\), and the assignments
\[
 e_1\longmapsto E_{11},\qquad
 e_2\longmapsto E_{22},\qquad
 a\longmapsto E_{21}
\]
give an isomorphism to the \(2\times 2\) lower triangular
matrix algebra over \(k\), with \(E_{ij}\) the usual matrix
units. Indeed, \(E_{22}E_{21}E_{11}=E_{21}\),
\(E_{21}^2=0\), the other vertex multiplications on \(E_{21}\)
are zero, and the diagonal matrix units are orthogonal and
sum to the identity. These are the products of the path basis.
The three matrix units span the triangular algebra and are linearly
independent. Here \(J^2=0\), so \(I=0\) is admissible.
Explicitly, a single map \(T\colon V_1\to V_2\) gives the action
\[
 \begin{pmatrix}x&0\\ u&y\end{pmatrix}(v_1,v_2)
   =\bigl(xv_1,\;uT(v_1)+yv_2\bigr).
\]
Indeed the displayed actions of \(E_{11},E_{22},E_{21}\) are the
two projections and the operator of \(a\) from \(V_1\) to \(V_2\).
\end{example}

\begin{example}[RM03-C06.10. One loop with a square-zero relation]
\label{ex:rm03-square-zero-loop}
For one vertex and one loop \(a\), the paths are \(a^n\)
for \(n\geq 0\), with \(a^0=e\) the vertex. The map
\(x^n\mapsto a^n\) is multiplicative on a basis and hence gives
\(k[x]\simeq kQ\). The relation \(a^2=0\) generates \(I=J^2\),
so it is admissible and \(kQ/I\simeq k[x]/(x^2)\) has basis
the classes of \(1,x\). A module is a vector space \(V\) with a
linear operator \(T=V_a\) satisfying \(T^2=0\). A morphism to
another such module with operator \(U\) is precisely a linear
map \(f\) with \(fT=Uf\). These claims follow from
Theorem~\ref{thm:rm03-bound-quiver-equivalence}, and the explicit
action is \((\lambda+\mu x)v=\lambda v+\mu T(v)\).
\end{example}


\backmatter
\begin{thebibliography}{9}

\bibitem{p103}
MathGenome, \emph{Algebra Core} (P1-03), as archived in MathGenome
v0.22.0, Chapters~1--3 and~8. Release source under
\texttt{provenance/components/AlgebraCore/delivery/}.

\bibitem{d07}
MathGenome, \emph{Abstract Algebra} (D07), revised lecture notes,
as archived in MathGenome v0.22.0, Chapters~9--11, 16, and~23.

\bibitem{lag}
MathGenome, \emph{Linear Algebraic Groups}, Chapter~6,
Results F01--F05. Companion lecture notes.

\bibitem{ktheory}
MathGenome, \emph{Algebraic K-Theory}, Chapter~1, Result A10.
Companion lecture notes.

\bibitem{mitdensity}
R. Bezrukavnikov, \emph{Noncommutative Algebra}, MIT 18.706, Spring 2023,
\href{https://ocw.mit.edu/courses/18-706-noncommutative-algebra-spring-2023/resources/mit18_706_s23_lec03_pdf/}{Lecture~3 on the density theorem}.

\bibitem{mitmorita}
R. Bezrukavnikov, \emph{Noncommutative Algebra}, MIT 18.706, Spring 2023,
\href{https://ocw.mit.edu/courses/18-706-noncommutative-algebra-spring-2023/resources/mit18_706_s23_lec07_pdf/}{Lecture~7 on Morita equivalence}.

\bibitem{stackscsa}
The Stacks Project Authors,
\href{https://stacks.math.columbia.edu/tag/073W}{\emph{Brauer groups}},
Chapter~11, Sections~11.3--11.8.

\bibitem{stacksdescent}
The Stacks Project Authors,
\href{https://stacks.math.columbia.edu/tag/03R1}{\emph{Brauer groups}},
Section~59.61, in \emph{\'Etale Cohomology}.

\bibitem{quiver}
W. Crawley-Boevey,
\href{https://www.math.uni-bielefeld.de/~wcrawley/25RA/Representations_of_Algebras_final_version_2025-07.pdf}{\emph{Representations of Algebras}},
July 2025, Sections~1.3--1.4.

\end{thebibliography}
\end{document}
